\section{Referee Proof Details: Terminal Object Before Estimate}

The proof begins with the terminal object because the Clay negation is a claim about one solution, not a claim about a failed method.
If smoothness fails at a finite time, the failure has to be carried by the same branch generated by the initial datum.
That branch has a material history, an equation history, and a field readout history.
The CM test separates those histories into three questions.
Pack asks whether there is still a positive same-fluid body.
Part asks whether the original Navier--Stokes equation is still the equation participating on that body.
Field asks whether that participating body still gives a positive-scale continuation readout.

This ordering is the proof's main discipline.
It prevents a bad picture from becoming a theorem by visual force.
A zero-radius remnant, an escaping high-frequency profile, a filtered pressure residue, or a point-like vortex picture may look like a singularity.
The proof asks a different question first: is that picture actually the terminal object of the same solution?
When it is not, it does not need to be smoothed.
It has already failed to be the in-class terminal branch required by the counterexample.

The proof also prevents a good estimate from being used on the wrong object.
A local \(H^s\) readout proves continuation only when the readout belongs to the same solution.
That is why Field is downstream of Pack and Part.
The argument cannot read a field window from a detached packet, a changed pressure law, or a rescaled surrogate and then declare the original solution continued.
The reader should therefore read every later estimate with this custody question in mind.

\section{Main Theorem}

\begin{theorem}[Whole-space same-solution terminal exclusion]
Let \(u_0\) be smooth, divergence-free, rapidly decaying data on \(\Rthree\), and let \(u\) be the corresponding classical Navier--Stokes solution with fixed viscosity \(\nu>0\).
There is no finite time \(T_*<\infty\) at which the same solution reaches an in-class nonsmooth terminal branch.
More precisely, every admitted terminal record \(Q\) for such a finite terminal claim satisfies exactly one of the following alternatives:
\[
\begin{aligned}
&\Pack_Q\wedge\Part_{N,Q}\wedge \exists r>0\,\Field_{N,r,Q}
&&\text{and the same solution continues,}\\
&\neg\Pack_Q
&&\text{and }Q\notin\Owork,\\
&\Pack_Q\wedge\neg\Part_{N,Q}
&&\text{and }Q\notin\Owork,\\
&\Pack_Q\wedge\Part_{N,Q}\wedge\forall r>0\,\neg\Field_{N,r,Q}
&&\text{and }Q\notin\Owork .
\end{aligned}
\]
The whole-space exterior source is included in the same alternatives: it either vanishes at the continuation readout level, or its dyadic survivor lands in Pack, Part, or Field.
\end{theorem}

\begin{proof}
Assume a finite terminal claim for the original solution.
The claim supplies a maximal interval \([0,T_*)\), the same smooth branch \(u\) on that interval, and a terminal record \(Q\) extracted from that branch along times tending to \(T_*\).
If no such record is supplied, the claim has not identified a terminal obstruction of the original solution.

Apply the ordered test.
If Pack fails, the record has no positive same-fluid body.
It is not a member of the working same-solution class.
If Pack survives and Part fails, the record has a body but the original pressure-viscosity equation no longer participates on that body.
It is again not a member of the working same-solution class.
If Pack and Part survive and every positive Field scale fails, the record has no continuation-level field readout and is not a member of the class.
These are exactly the three face failures defining \(\Exit(Q;\Owork)=\neg\Member(Q;\Owork)\).

The only remaining case is the pass branch.
There is a positive same-fluid packet, the original equation participates on it, and a positive-scale \(H^s\), \(s>5/2\), readout is retained.
The standard local continuation theorem restarts the same branch past \(T_*\).
That contradicts terminality.

It remains to account for the whole-space exterior source.
The Duhamel decomposition separates the linear heat tail, compact-core nonlinear source, annular bookkeeping, and genuinely exterior nonlinear source.
The first three do not create an independent finite terminal high-order exterior branch.
If the genuinely exterior source vanishes in the exterior \(H^s\) readout, the pass-side input is supplied.
If it does not vanish, bounded exterior frequencies are removed by exterior \(L^2\) tightness, so any survivor occurs along dyadic frequencies \(N_j\to\infty\).
Such a survivor either lacks a same-fluid material address, lacks same-equation participation, or contradicts any fixed positive Field readout.
Thus the exterior survivor also lands in Pack, Part, or Field.
No third terminal branch remains.
\end{proof}

\section{Terminal Record Entry In Standard PDE Language}

The terminal-record step does not rely on private vocabulary.
It is the ordinary maximal-solution argument written in the language needed for the proof.
Let \(s>5/2\).
Local well-posedness says that an \(H^s\) datum generates a unique classical branch for a positive time depending on the \(H^s\) size of the datum.
It also says that a classical branch can be restarted from any time at which the \(H^s\) size is finite.

\begin{lemma}[Finite terminality produces a same-branch readout failure]
Assume \(u\) is a maximal classical solution on \([0,T_*)\) with \(T_*<\infty\) and no classical continuation past \(T_*\).
Then for every continuation exponent \(s>5/2\),
\[
  \limsup_{t\uparrow T_*}\|u(t)\|_{H^s}=\infty .
\]
\end{lemma}

\begin{proof}
If \(\|u(t)\|_{H^s}\) remained bounded along the interval, the local theory would give a restart time bounded below uniformly along times \(t\) close to \(T_*\).
Choosing \(t<T_*\) close enough would extend the same solution past \(T_*\).
That contradicts maximality.
\end{proof}

The proof does not need a strong endpoint limit.
It needs a same-branch terminal record.
The record can be a sequence \(t_j\uparrow T_*\), together with the material packets, equation participation data, and field readouts being tested along that sequence.
The record is not assumed to be a class member.
Membership is the result of the test, not the entry condition.

\begin{lemma}[Entry does not assume membership]
Let \(Q\) be the record extracted from a finite terminal claim.
The statement that \(Q\) enters the CM test says only that \(Q\) is offered as a terminal object of the original solution.
It does not assert \(\Member(Q;\Owork)\).
\end{lemma}

\begin{proof}
Entry fixes the referent.
The record comes from the same datum, same branch, same viscosity, and same pressure law before the terminal test is applied.
After entry, the ordered faces decide whether the record remains a member.
Therefore entry and membership are logically different.
\end{proof}

\section{The Working Class Is A Same-Solution Class}

The working class \(\Owork\) is deliberately narrower than the set of all patterns one might see near a possible singular time.
It contains only records that still carry the three structures needed to be a terminal branch of the original solution.
Those structures are material address, equation participation, and continuation-level field readout.

\begin{definition}[Membership]
For the purposes of the proof,
\[
  \Member(Q;\Owork)
\]
means that \(Q\) retains Pack, Part, and at least one positive Field scale at the continuation depth being used.
The ordered form is
\[
  \Pack_Q\wedge\Part_{N,Q}\wedge \exists r>0\,\Field_{N,r,Q}.
\]
\end{definition}

\begin{proposition}[Ordered negation]
For an admitted record \(Q\), failure of membership is detected by the first failed ordered face:
\[
\neg\Member(Q;\Owork)
\]
is represented by
\[
\neg\Pack_Q,
\quad
\Pack_Q\wedge\neg\Part_{N,Q},
\quad\text{or}\quad
\Pack_Q\wedge\Part_{N,Q}\wedge\forall r>0\,\neg\Field_{N,r,Q}.
\]
\end{proposition}

\begin{proof}
Membership is the conjunction of the three ordered requirements.
The negation of an ordered conjunction is exactly the first place where the conjunction fails.
The ordered form matters because Part is meaningful only on a packed material carrier, and Field is meaningful only on a participating carrier.
\end{proof}

This is not a trick definition.
The Clay terminal claim needs the same structures.
Without a material address, there is no terminal body of the original fluid.
Without equation participation, there is no Navier--Stokes terminal branch of the original equation.
Without a positive field readout, there is no continuation-side member branch.

\section{Pack Payment: Material Body Of The Same Fluid}

Pack is the first face because a terminal obstruction needs a body before it can be acted on by the equation or read as a field.
At a finite terminal time, many bad scenarios can be described only after passing to rescaled variables, taking weak limits, or focusing on a formal point.
Those scenarios may be useful for analysis.
They are not automatically the same solution's terminal body.

\begin{definition}[Positive material body]
A positive material body for \(Q\) is a packet with positive physical scale whose ancestry is transported by the original velocity field on preterminal times.
The packet is allowed to be selected by the terminal argument.
It is not allowed to exist only after changing the problem.
\end{definition}

\begin{lemma}[Pack excludes detached terminal pictures]
If a proposed obstruction has zero physical radius, no same-fluid ancestry, or exists only as a rescaled profile detached from the original variables, then it fails Pack.
\end{lemma}

\begin{proof}
Pack asks for a positive same-fluid packet in the original solution.
A zero-radius object has no positive carrier.
A profile with no transported ancestry is not the original fluid packet.
A rescaled diagnostic may describe a limiting shape, while Pack keeps the decisive question on the positive packet in the original solution before Part and Field are asked.
Thus each detached picture gives \(\neg\Pack_Q\).
\end{proof}

\begin{corollary}[Pack failure is not a hidden smoothness claim]
When Pack fails, the proof has not shown that the bad picture is smooth.
It has shown that the picture is not the same solution's in-class terminal member.
\end{corollary}

\section{Part Payment: The Same Equation Participates}

Part is the second face.
It is needed because a localized packet can be visually tied to the same fluid while the equation being used on it has changed.
Cutoffs introduce commutators.
Pressure is nonlocal.
Periodic replacements change spatial infinity.
Filters and comparison branches can preserve some estimates while changing the object.
Part prevents those changes from being smuggled into the proof.

\begin{definition}[Same-equation participation]
\(\Part_{N,Q}\) means that through depth \(N\), the packet-level evolution used in the proof is inherited from
\[
  \partial_t u+\mathbb P\nabla\cdot(u\otimes u)=\nu\Delta u
\]
for the original whole-space branch, with the pressure, viscosity, nonlinear source, commutators, and cutoffs accounted for as terms of that branch.
\end{definition}

\begin{lemma}[Changed equation gives Part failure]
Assume Pack survives.
If the packet is acted on by a filtered equation, a comparison flow, a periodic substitute, a discarded pressure law, or an unaccounted cutoff commutator, then Part fails until the changed term is returned to the same-equation accounting.
\end{lemma}

\begin{proof}
Part requires that the original equation participates on the selected packet.
The listed objects may be valid support tools, but they do not by themselves prove participation of the original equation.
If the difference term is not carried back into the packet equation, the packet is no longer being tested as the original Navier--Stokes branch.
\end{proof}

This is the reason the whole-space proof is not a periodic proof in disguise.
The periodic material atlas gives a clean terminal-capture model without spatial infinity.
The whole-space proof still has to carry the \(\Rthree\) pressure and exterior source through Part.
It does so through the Duhamel split and the exterior-source analysis below.

\section{Field Payment: Positive-Scale Continuation Readout}

Field is the third face.
It is the place where the proof asks for a continuation-level readout after material and equation custody have survived.
The readout is not a global all-scales miracle.
It is a positive-scale field window at a Sobolev depth high enough for restart.

\begin{definition}[Positive Field scale]
\(\Field_{N,r,Q}\) means that on a positive physical scale \(r>0\), after Pack and Part have survived, the selected packet carries coherent field data through depth \(N\) strong enough to place the same branch in the \(H^s\), \(s>5/2\), continuation framework.
\end{definition}

\begin{lemma}[Field pass gives restart]
If \(\Pack_Q\), \(\Part_{N,Q}\), and \(\Field_{N,r,Q}\) hold at a continuation depth, then the selected record supplies restart data for the same solution.
\end{lemma}

\begin{proof}
Pack gives the same-fluid packet.
Part gives the original equation on that packet.
Field gives a positive-scale \(H^s\) readout.
The usual continuation theorem restarts the same solution from the readout.
\end{proof}

\begin{lemma}[Field failure is the terminal high-order failure]
Assume Pack and Part survive.
If no positive physical scale carries the continuation-level readout, then the record fails Field and exits the working class.
\end{lemma}

\begin{proof}
After Pack and Part, the record is still the same solution's participating packet.
The only remaining question for membership is whether a positive continuation-scale field readout exists.
If no such scale exists, \(\exists r>0\,\Field_{N,r,Q}\) is false, so \(Q\notin\Owork\).
\end{proof}

\section{Pass Branch And Fail Branch}

The proof has two outcomes.
The pass branch is a genuine continuation branch.
The fail branch is a class-exit branch.
They are not two ways for a Clay counterexample to survive.

On the pass branch,
\[
  \Pack_Q\wedge\Part_{N,Q}\wedge\exists r>0\,\Field_{N,r,Q}
\]
identifies an in-class same-solution record.
The continuation readout is then attached to the original solution.
The standard local theorem extends that solution.
Thus a terminal counterexample cannot lie on the pass branch.

On the fail branch, the terminal record has lost one of the structures needed to be the original solution's in-class terminal member.
It is not a member branch with nonsmooth terminal behavior.
It is the first face at which the alleged terminal object stops being a legal terminal member.
Thus a terminal counterexample cannot lie on the fail branch either, unless it can produce a fourth kind of same-solution terminal object that defeats the three membership faces.
The later sections show that the whole-space exterior source does not produce such a fourth object.

\section{Periodic Atlas As The Compact Model}

On \(\Tthree\), there is no spatial infinity.
The clean terminal model is a finite transported material atlas.
Choose a finite smooth atlas and partition of unity on the torus.
Transport it by the same velocity field on preterminal times.
The terminal record is the terminal tail of the whole transported atlas.
This object covers the entire periodic solution, so it cannot miss a terminal branch by choosing the wrong local packet.

\begin{proposition}[Periodic terminal atlas exhaustion]
For a finite periodic terminal claim, the transported atlas either loses Pack, loses Part after Pack, loses every positive Field scale after Pack and Part, or keeps all three and continues.
\end{proposition}

\begin{proof}
The finite atlas supplies a finite collection of same-fluid packets covering the torus.
If the transported atlas loses positive material compatibility, Pack fails.
If it keeps material compatibility but the original pressure-viscosity equation no longer participates through the required depth, Part fails.
If Pack and Part survive but no positive local field window remains, Field fails.
If Pack, Part, and one positive Field scale survive, the finite overlap of the atlas sums the local readouts into the global \(H^s\) continuation norm, and the standard continuation theorem extends the solution.
\end{proof}

This periodic model explains the logic of the faces.
The whole-space proof uses the same membership logic, with additional work to keep spatial infinity from becoming a fourth branch.

\section{Whole-Space Duhamel Decomposition}

In \(\Rthree\), the exterior region has to be treated before the terminal packet test is complete.
Let \(\eta_R\) be an exterior cutoff, equal to one outside a large radius and supported farther out.
Write the Leray-Duhamel formula schematically as
\[
u(t)=e^{\nu t\Delta}u_0
-\int_0^t e^{\nu(t-\tau)\Delta}\mathbb P\nabla\cdot(u\otimes u)(\tau)\,d\tau.
\]
Split the nonlinear source into a compact-core part, an annular bookkeeping part, and a genuinely exterior part.
The exterior contribution is then decomposed as
\[
  \eta_R u(t)=L_R(t)+C_R(t)+A_R(t)+W_R(t),
\]
where \(L_R\) is the exterior heat tail of the initial data, \(C_R\) is the far-field response of the compact-core source, \(A_R\) contains cutoff and annular terms, and \(W_R\) is the exterior nonlinear source.

\begin{lemma}[Initial tail removal]
For smooth rapidly decaying \(u_0\), the heat tail \(L_R\) vanishes in \(H^s\) as \(R\to\infty\) on every finite preterminal interval.
\end{lemma}

\begin{proof}
Rapid decay gives weighted Sobolev control of all orders.
The heat semigroup preserves and improves this decay on finite time intervals.
After differentiating up to order \(s\), the exterior weighted tail is bounded by a negative power of \(R\).
\end{proof}

\begin{lemma}[Compact-core far-kernel removal]
For a fixed compact source region, the far-field Duhamel response \(C_R\) vanishes in the exterior \(H^s\) readout as \(R\to\infty\).
\end{lemma}

\begin{proof}
The heat kernel and Leray-projected derivative kernel have off-diagonal decay away from the source support.
The compact source is separated from the exterior cutoff by a distance comparable to \(R\).
The finite number of derivatives needed for the \(H^s\) readout is controlled on the compact preterminal region.
The off-diagonal kernel bounds therefore drive the exterior contribution to zero.
\end{proof}

\begin{lemma}[Annular bookkeeping has no independent terminal role]
The annular terms either vanish with the cutoff bookkeeping or are absorbed into the same exterior source record \(W_R\).
\end{lemma}

\begin{proof}
Annular terms are generated only by the chosen partition.
They are not a separate equation and not a separate solution branch.
If the terms are small, they are harmless.
If they persist, their persistence is a part of the exterior-source packet under test.
\end{proof}

After these reductions, the only possible exterior obstruction is \(W_R\).
The whole-space proof therefore needs to decide whether \(W_R\) vanishes at the continuation readout level or survives as a terminal exterior packet.

\section{Exterior Source Alternative}

The exterior nonlinear source has exactly two alternatives:
\[
  \lim_{R\to\infty}\sup_{t<T_*}\|W_R(t)\|_{H^s}=0,
\]
or there are \(R_j\to\infty\) and \(t_j<T_*\) for which
\[
  \|W_{R_j}(t_j)\|_{H^s}\ge\varepsilon>0.
\]
The first alternative supplies the positive exterior input.
The second alternative has to be classified.

\begin{proposition}[Bounded frequencies cannot carry the surviving exterior branch]
Assume the exterior \(L^2\) tightness of the original solution and the Duhamel split above.
For every fixed dyadic cutoff \(N_0\),
\[
  \lim_{R\to\infty}\sup_{t<T_*}
  \|P_{\le N_0}W_R(t)\|_{H^s}=0.
\]
\end{proposition}

\begin{proof}
On bounded frequency,
\[
  \|P_{\le N_0}f\|_{H^s}\le C_sN_0^s\|f\|_{L^2}.
\]
The exterior \(L^2\) contribution of the original solution is tight.
The initial and compact-core parts have already been removed.
The bounded-frequency exterior response therefore tends to zero as \(R\to\infty\).
\end{proof}

\begin{corollary}[Dyadic survivor]
If the exterior nonlinear source does not vanish in \(H^s\), then there exist \(R_j\to\infty\), \(t_j<T_*\), and dyadic frequencies \(N_j\to\infty\) such that
\[
  N_j^s\|P_{N_j}W_{R_j}(t_j)\|_{L^2}\ge c>0.
\]
\end{corollary}

\begin{proof}
If every high dyadic frequency vanished as well, the \(H^s\) norm would vanish by summing the bounded and high-frequency parts.
The positive lower bound therefore has to move to dyadic frequencies escaping to infinity.
\end{proof}

\section{Dyadic Survivor Landing}

The dyadic survivor is not a fourth terminal branch.
It is a high-frequency exterior packet.
The CM test asks whether that packet has a same-fluid material address, whether the same equation participates on it, and whether it can coexist with a fixed positive Field scale.

\begin{proposition}[Dyadic survivor lands in Pack, Part, or Field]
Let \(Q\) be the terminal record selected by a dyadic exterior survivor.
Then the survivor forces one of
\[
  \neg\Pack_Q,\qquad
  \Pack_Q\wedge\neg\Part_{M,Q},\qquad
  \Pack_Q\wedge\Part_{M,Q}\wedge\forall\rho>0\,\neg\Field_{M,\rho,Q}
\]
for continuation depth \(M>s\).
\end{proposition}

\begin{proof}
If the high-frequency survivor has no retained material address in the original variables, then it has no positive same-fluid packet and Pack fails.
Assume Pack survives.
If the survivor is not acted on by the original Navier--Stokes equation with the original pressure, viscosity, and nonlinear source ancestry, then Part fails.
Assume Pack and Part survive.
The survivor is now a same-solution participating packet.
Suppose, toward contradiction, that \(\Field_{M,\rho,Q}\) holds for some fixed \(\rho>0\).
On each retained \(\rho\)-scale Field window, finite derivative control through depth \(M\) gives
\[
  \|P_N f\|_{L^2}\le C_{\rho,M,Q}N^{-M}.
\]
Multiplying by \(N^s\) gives
\[
  N^s\|P_N f\|_{L^2}\le C_{\rho,M,Q}N^{s-M}\to0.
\]
The finite number of retained Field windows gives the same decay for the selected packet.
This contradicts the dyadic lower bound along \(N_j\to\infty\).
Therefore no fixed positive Field scale survives.
\end{proof}

The proof uses the Field assumption only to reject the Field branch.
It does not assume global \(H^s\) control.
It says that a fixed positive Field window cannot be the same object as a survivor whose \(H^s\) mass escapes to arbitrarily high dyadic frequency.

\section{Zero-Radius And Vortex-Like Obstructions}

A common endpoint picture is shrinking scale.
The proof treats it directly.
If the alleged obstruction has collapsed to radius zero, it has no positive same-fluid packet.
That is Pack failure.
The proof does not need to show that the zero-radius picture is regular.
It needs to show that it is not a positive-body terminal member of the original branch.

A vortex-like picture requires the same care.
A vortex is not merely high rotation at a point.
As a fluid object, it has spatial extent, circulation structure, material ancestry, and field readout.
A point-like vortex with no positive radius has already lost the object it claims to be.
If a positive vortex packet survives but its pressure-viscosity equation has been replaced by a local model or rescaled limit, Part fails until the original equation is restored.
If the same packet and equation survive but every positive scale loses coherent continuation readout, Field fails.
Thus the vortex picture is handled by the same ordered test rather than by an ad hoc vortex theorem.

\section{No Hidden Positive Estimate}

The proof does not claim to have proved a direct noncircular global \(H^s\) bound.
It proves a pass-or-exit theorem for admitted terminal records.
This distinction is essential.
A direct exterior high-order estimate would try to show
\[
  \lim_{R\to\infty}\sup_{t<T_*}\|\eta_Ru(t)\|_{H^s}=0
\]
by estimating all exterior high derivatives.
That route becomes circular if it spends a continuation-grade coefficient such as a global \(H^s\) bound or a Beale--Kato--Majda integral.

The CM route avoids that circularity.
When the exterior tail vanishes, the pass input is available.
When the exterior tail survives, the proof does not try to smooth it by force.
It classifies the survivor as a terminal packet.
The survivor then has to keep Pack, Part, and Field.
If it keeps them, it cannot remain a dyadic high-frequency survivor.
If it loses one, it exits the class.

\section{Referee Objections In The Proof Language}

\subsection*{Objection: class exit is just another name for blowup}

Class exit is not a name for blowup.
It is the conclusion that a tested terminal record has lost a required same-solution membership face.
The alleged counterexample needs an in-class nonsmooth terminal branch.
The fail branch supplies a reason that the alleged record is not that branch.

\subsection*{Objection: the fail branch could still be the Clay obstruction}

The Clay obstruction is an obstruction to continuing the original solution.
If the record fails Pack, it lacks a positive body of that solution.
If it fails Part, it lacks the same equation on that body.
If it fails Field after Pack and Part, it lacks the continuation readout needed for membership.
The fail branch is therefore the place where the alleged terminal object stops satisfying the same-solution terminal-member requirements.
It is not a hidden member-side singularity.

\subsection*{Objection: the pass branch could be a different branch}

The pass branch is downstream of Pack and Part.
Pack ties the readout to a same-fluid packet.
Part ties the packet to the original pressure-viscosity equation.
Only then is Field allowed to call the continuation theorem.
That ordering is exactly what prevents the pass branch from becoming a different branch.

\subsection*{Objection: the whole-space exterior source was not proved small}

The proof does not require the exterior source to be small in every case.
It requires a complete decision.
If the exterior source is small at \(H^s\) scale, it supplies the pass input.
If it is not small, bounded frequencies cannot carry it, so it has a dyadic survivor.
The dyadic survivor is then tested by Pack, Part, and Field.
That is the whole-space decision.

\section{Final Assembly Of The Contradiction}

Assume a finite Clay terminal counterexample.
The counterexample gives the same datum, same viscosity, same pressure law, same preterminal branch, and a terminal record \(Q\).
The record enters the ordered CM test.

If all faces survive, the record is a same-solution member branch with a continuation-level Field readout.
The standard \(H^s\) theorem restarts the same solution past the alleged terminal time.
This contradicts terminality.

If a face fails, the first failed face gives \(\Exit(Q;\Owork)=\neg\Member(Q;\Owork)\).
The alleged obstruction is not an in-class nonsmooth terminal branch of the original solution.
The whole-space exterior source cannot avoid this split: it either supplies the pass input or becomes a dyadic survivor, and the dyadic survivor lands in the same face test.

Thus the terminal counterexample has no remaining branch.
The pass branch continues.
The fail branch exits the same-solution member class.
There is no third whole-space exterior branch.
Therefore the original smooth rapidly decaying whole-space Navier--Stokes solution has no finite same-solution terminal nonsmooth member branch.
